\slajdid{04-s008}
\begin{frame}[label=04-s008]{Celobrojni deo: uzastopno deljenje}
% element: 04-s008:e01 — Deljenje CV10 sa r i izdvajanje ostatka c0, uz c0/r<1 i skup cifara
% element: 04-s008:e02 — Naredno deljenje količnika i ostatak c1
% element: 04-s008:e03 — Uslov zaustavljanja i vodeća cifra cn−1; nastavak daje vodeće nule
\setlength{\abovedisplayskip}{3pt}\setlength{\belowdisplayskip}{3pt}
Pri deljenju celobrojnog dela sa $r$:
\[\begin{aligned}CV_{10}\div r&=c_{n-1}r^{n-2}+c_{n-2}r^{n-3}+\ldots+c_1r^0+\frac{c_0}{r}\\&=(CV_{10})_1+\frac{c_0}{r}\end{aligned}\]
\[\frac{c_0}{r}<1,\qquad c_i\in\{0,1,\ldots,r-1\}\]
Izdvajamo ostatak $\alert{c_0}$ i nastavljamo sa količnikom:
\[(CV_{10})_1\div r=(CV_{10})_2+\frac{c_1}{r}.\]
Izdvajamo $\alert{c_1}$ i ponavljamo postupak.
\Rezultat{Zaustavljanje: $(CV_{10})_{n-1}<r$, pa je $(CV_{10})_{n-1}=c_{n-1}$.}
\[(CV_{10})_n=0\qquad\text{Dalje bismo dobijali vodeće nule.}\]
\end{frame}
